\section{Resultados esperados y productos I+D+i}

En el marco del desarrollo del presente proyecto, se esperan obtener resultados en 
Ciencia, Tecnología e Innovación (CTeI) alineados con las tipologías establecidas por 
MinCiencias, los cuales se describen a continuación:

\subsection{Productos de generación de nuevo conocimiento}
\begin{itemize}
    \item Un (1) artículo científico sometido a una revista indexada en categoría Q1 
    o Q2, relacionado con el modelado computacional (CFD), control y optimización del 
    proceso de secado de café basado en humedad relativa mediante bomba de calor.
    \item Un (1) documento técnico o capítulo de libro que aborde el análisis del 
    proceso de secado y su impacto en la calidad física, química y sensorial de cafés 
    diferenciados (\textit{honey} y \textit{natural}).
\end{itemize}

\subsection{Productos de desarrollo tecnológico e innovación}
\begin{itemize}
    \item Diseño y validación de un sistema de secado de café a escala piloto avanzado, 
    basado en bomba de calor y control de humedad relativa, con criterios de diseño 
    que incorporan consideraciones de escalabilidad para su futura transferencia 
    tecnológica al sector productivo.
    \item Prototipo funcional a escala piloto avanzado con estrategias de manejo del 
    producto —sistemas de movimiento y agitación controlada— orientadas a garantizar 
    la homogeneidad del secado en procesos de cafés diferenciados, acompañado de las 
    especificaciones técnicas necesarias para su replicabilidad.
\end{itemize}

\subsection{Productos de apropiación social del conocimiento}
\begin{itemize}
    \item Realización de al menos un (1) taller de socialización con actores del sector 
    cafetero (caficultores, técnicos y entidades aliadas), orientado a la transferencia 
    de conocimiento sobre tecnologías de secado controlado para cafés diferenciados.
    \item Elaboración de una guía o protocolo de operación del sistema de secado, 
    dirigida a operarios y técnicos del sector productivo, que documente de forma 
    clara y replicable las condiciones óptimas de operación identificadas para procesos 
    \textit{honey} y \textit{natural}.
    \item Producción de contenidos de divulgación en formatos no académicos —infografías, 
    videos y podcasts— dirigidos al público general y a las comunidades cafeteras del 
    Valle del Cauca, con el propósito de democratizar el conocimiento generado y 
    fortalecer la apropiación social de la tecnología desarrollada.
\end{itemize}

\subsection{Productos de formación de recurso humano en CTeI}
\begin{itemize}
    \item Formación de un (1) estudiante de Maestría en Ingeniería Agroindustrial 
    mediante el desarrollo del presente proyecto de investigación.
    \item Vinculación de un (1) estudiante de pregrado en Ingeniería Agroindustrial 
    mediante la modalidad de trabajo de grado, participando en actividades 
    experimentales y de análisis de resultados en el marco del proyecto.
\end{itemize}

Los resultados propuestos son coherentes con el alcance de la investigación a nivel 
piloto avanzado y buscan contribuir tanto al avance del conocimiento científico como 
al desarrollo de soluciones tecnológicas aplicables al sector cafetero colombiano. 
Si bien la implementación a escala industrial completa trasciende el tiempo de 
ejecución de la maestría, los productos generados constituirán la base técnica, 
experimental e institucional para futuros procesos de transferencia tecnológica, 
escalamiento industrial y fortalecimiento de la agroindustria cafetera del Valle 
del Cauca.